%mac
\loai{Ví dụ mẫu}
\begin{vidu}	
Quan sát chuyển động của xe ở hình vẽ sau?
\begin{center}
\includegraphics[width=0.6\linewidth]{images/8-1}
\end{center}
\begin{itemize}
\item[a)] Tính gia tốc của ô tô trên $4$ đoạn đường trong hình.
\item[b)] Gia tốc của ô tô trên đoạn đường $4$ có gì đặc biệt so với gia tốc trên các đoạn đường khác?
\end{itemize}	
\textbf{Hướng dẫn:}
\begin{itemize}
\item[a)]  Đổi $5\ km/h=.............m/s;\ 29\ km/h=.............m/s;\ 49\ km/h=.............m/s;\ 30\ km/h=\dotfill m/s$.
\item[-] Gia tốc trong đoạn đường 1: $a_1=\dfrac{\Delta v_1}{\Delta t}=$\dotfill 
\item[-] Gia tốc trong đoạn đường 2: $a_2=$\dotfill 
\item[-] Gia tốc trong đoạn đường 3: $a_3=$\dotfill 
\item[-] Gia tốc trong đoạn đường 4: $a_4=$\dotfill 
\item[b)] Đoạn đường 1, 2 và 3 gia tốc\dotfill \\
Đoạn đường 4 gia tốc\dotfill 
\end{itemize}
\loigiai{
\begin{itemize}
\item[a)]  Đổi $5\ km/h=\dfrac{25}{18}\ m/s;\ 29\ km/h=\dfrac{145}{18}\ m/s;\ 49\ km/h=\dfrac{245}{18}\ m/s;\ 30\ km/h=\dfrac{15}{3}\ m/s$.
\item[-] Gia tốc trong đoạn đường 1: $a=\dfrac{\Delta v}{\Delta t}=\dfrac{25}{18.1}=\dfrac{25}{18}\approx 1,39\ m/s^2$
\item[-] Gia tốc trong đoạn đường 2: $a=\dfrac{\Delta v}{\Delta t}=\dfrac{\dfrac{145}{18}-\dfrac{25}{18}}{4-1}\approx 2,22\ m/s^2$
\item[-] Gia tốc trong đoạn đường 3: $a=\dfrac{\Delta v}{\Delta t}=\dfrac{\dfrac{245}{18}-\dfrac{145}{18}}{6-4}\approx 2,78\ m/s^2$
\item[-] Gia tốc trong đoạn đường 4: $a=\dfrac{\Delta v}{\Delta t}=\dfrac{\dfrac{25}{3}-\dfrac{245}{18}}{7-6}\approx -5,28\ m/s^2$
\item[b)] Trong 4 đoạn đường trên, vận tốc tăng dần, còn gia tốc từ đoạn đường 1 đến đoạn đường 3 tăng dần, nhưng từ đoạn đường 3 đến đoạn đường 4 thì gia tốc giảm dần.
\end{itemize}
}
\haidong{5}
\end{vidu}
\begin{vidu}
Một xe máy đang chuyển động thẳng với vận tốc $10\ m/s$ thì tăng tốc. Sau $5\ s$ đạt vận tốc $12\ m/s$.
\begin{vdc}
Gia tốc của xe là bao nhiêu $m/s^2$?
\shortans[oly]{0,4}
\loigiai{
\begin{itemize}
\item $v_0=10\ m/s;\ v=12\ m/s;\ \Delta t=5\ s$.
\item $a=\dfrac{\Delta v}{\Delta t}=\dfrac{12-10}{5}=0,4\ m/s^2$.
\item Gia tốc của xe là $a=0,4\ m/s^2$.
\end{itemize}	
}	
\end{vdc}
\begin{vdc}
Nếu sau khi đạt vận tốc $12\ m/s$, xe chuyển động chậm dần với gia tốc có độ lớn bằng gia tốc trên thì sau bao nhiêu giây xe sẽ dừng lại?
\shortans[oly]{30}
\loigiai{
\begin{itemize}
\item $v_0'=12\ m/s;\ v'=0;\ a=-0,4\ m/s^2$.
\item $\Delta t'=\dfrac{\Delta v'}{a}=\dfrac{0-12}{-0,4}=30\ s$.
\item Xe dừng lại sau 30 s.
\end{itemize}	
}	
\end{vdc}
\end{vidu}
\begin{vidu}
Một vận động viên đua xe đạp đường dài vượt qua vạch đích với tốc độ $10\ m/s$. Sau đó vận động viên này đi chậm dần đều thêm $20\ m$ nữa mới dừng lại. Coi chuyển động của vận động viên là thẳng.
\begin{vdc}
Gia tốc của vận động viên trong đoạn đường sau khi đi qua vạch đích có độ lớn là bao nhiêu $m/s^2$?
\shortans[oly]{2,5}
\loigiai{
\begin{itemize}
\item Gia tốc của vận động viên trong đoạn đường sau khi qua vạch đích là: 
\item $v^2-v_0^2=2ad\Rightarrow a=\dfrac{v^2-v_0^2}{2d}=\dfrac{0^2-10^2}{2.20}=-2,5\ m/s^2$.
\end{itemize}
}
\end{vdc}
\begin{vdc}
Thời gian vận động viên đó cần để dừng lại kể từ khi cán đích là bao nhiêu giây?
\shortans[oly]{4}
\end{vdc}
\begin{vdc}
Vận tốc trung bình của người đó trong quá trình dừng xe là bao nhiêu $m/s$?
\shortans[oly]{5}
\end{vdc}
\loigiai{
\begin{itemize}
\item Vận tốc trung bình của người đó trên quãng đường dừng xe là: $v=\dfrac{d}{t}=\dfrac{20}{4}=5\ m/s$.
\end{itemize}
}
\end{vidu}





\begin{vidu}%SBT-KNTT9.8
Một ô tô khi hãm phanh có thể có gia tốc tối đa có độ lớn là $3\ m/s^2$. Ô tô đang chạy với vận tốc là $72\ km/h$. 
\begin{vdc}
Ô tô phải hãm phanh cách vật cản là bao nhiêu mét để không đâm vào vật cản (làm tròn kết quả đến chữ số hàng phần mười)?  
\shortans[oly]{66,7}
\loigiai{
\begin{itemize}
\item $v^2-v_0^2=2as\Rightarrow d=\dfrac{v^2-v_0^2}{2a}=\dfrac{0^2-20^2}{2.(-3)}\approx 66,7\ m$.
\end{itemize}
}
\end{vdc}
\begin{vdc}
Thời gian hãm phanh là bao nhiêu giây (làm tròn kết quả đến chữ số hàng phần mười)?
\shortans[oly]{6,7}
\loigiai{
\begin{itemize}
\item $v=v_0+at\Rightarrow t=\dfrac{v-v_0}{a}=\dfrac{-20}{-3}\approx 6,7\ s$.
\end{itemize}
}
\end{vdc}
\end{vidu}

\begin{vidu}%[1P1-6.1K][Loai1]
Một vật chuyển động thẳng biến đổi đều có phương trình vận tốc theo thời gian $t$ là: $v_t = 20 + 8t$ (trong đó $x$ đo bằng $cm$, $t$ đo bằng $s$).
\begin{vdc}
Vận tốc đầu của vật là bao nhiêu $m/s$?
\shortans[oly]{20}
\loigiai{
So sánh với phương trình vận tốc tổng quát $v_t = v_0 + at \Rightarrow  v_0=20\ cm/s$
}
\end{vdc}
\begin{vdc}
Gia tốc của vật là bao nhiêu $m/s^2$?
\shortans[oly]{8}
\loigiai{
So sánh với phương trình vận tốc tổng quát $v_t = v_0 + at \Rightarrow a=8\ cm/s^2$
}
\end{vdc}

\begin{vdc}
Quãng đường đi được trong khoảng thời gian $t = 5\ s$ kể từ khi bắt đầu chuyển động là bao nhiêu $cm$?
\shortans[oly]{200}
\loigiai{
Vận tốc của vật khi $t = 5$s là: $v_t = 20 + 8.5 = 60$ cm/s
\begin{itemize}
\item Quãng đường đi được của vật khi $t = 5$ s: $s=\dfrac{v_t^2-v_0^2}{2a}=\dfrac{60^2-20^2}{2.8}=200\ cm$
\end{itemize}
}
\end{vdc}
\begin{vdc}
Quãng đường vật đi được từ $t_1 = 2\ s$ đến $t_2 = 5\ s$ là bao nhiêu $cm$?
\shortans[oly]{144}
\loigiai{
Vận tốc của vật khi $t_1 = 2$s: $v_1 = 20 + 8.2 = 36$ cm/s
\begin{itemize}
\item Quãng đường vật đi được trong thời gian $t_1 = 2$s: 	$s_1=\dfrac{v_1^2-v_0^2}{2a}=\dfrac{36^2-20^2}{2.8}=56\ cm$
\item Vận tốc của vật khi $t_2 = 5$s: $v_2 = 20 + 8.5 = 60$ cm/s
\item Quãng đường vật đi được trong thời gian $t_2 = 5$ s: 	$s_2=\dfrac{v_2^2-v_0^2}{2a}=\dfrac{60^2-20^2}{2.8}=200\ cm$
\item Vậy quãng đường chất điểm đi được trong thời gian từ $t_1 = 2$ s đến $t_2 = 5$ s
\begin{align*}
\Delta s=s_2-s_1=200-56=144\ cm
\end{align*}
\end{itemize}}
\end{vdc}
\end{vidu}
\begin{vidu}%[1P1-5.3K][Loai1] 
Viên bi chuyển động nhanh dần đều không vận tốc đầu với gia tốc $a = 0,2\ m/s^2$.
\begin{vdc}
Quãng đường viên bi đi được trong $6$ giây là bao nhiêu mét?
\shortans[oly]{3,6}
\loigiai{
\begin{itemize}
\item Ta có: $s=\dfrac{1}{2}at^2=\dfrac{0,2}{2}6^2=3,6\ m$.
\end{itemize}	
}
\end{vdc}
\begin{vdc}
Quãng đường viên bi đi được trong giây thứ $6$ là bao nhiêu mét?
\shortans[oly]{1,1}
\loigiai{
\begin{itemize}
\item Ta có: $s=\dfrac{1}{2}at^2\Rightarrow \Delta s=s_{(t=6)}-s_{(t=5)}=\dfrac{a}{2}\left(6^2-5^2\right)=1,1\ m$.
\end{itemize}	
}
\end{vdc}
\end{vidu}

\loai{Chuyển động nhanh dần đều}
\begin{vidu}%[1P1B5-2][Loai1]
Một xe lửa bắt đầu rời khỏi ga và chuyển động thẳng nhanh dần đều với gia tốc $0,1\ m/s^2$. Khoảng thời gian để xe lửa đạt được vận tốc $36\ km/h$ là
\choice
{$360\ s$}
{\True $100\ s$}
{$300\ s$}
{$200\ s$}	
\haidong{14}
\loigiai{
\begin{itemize}
\item Đổi $36\ km/h = 10\ m/s$.
\item Ta có: $v_t=v_0+at \Rightarrow t=\dfrac{v_t-v_0}{a}=\dfrac{10}{0,1}=100\ s$. 
\end{itemize}	
}
\end{vidu}

\begin{vidu}%[1P1K5-2][Loai1]
Một xe đang chuyển động với vận tốc $18\ km/h$ thì tăng tốc đều với gia tốc $a$. Sau $100\ m$ xe đạt vận tốc $54\ km/h$. Giá trị gia tốc a là
\choice
{$0,5\ m/s^2$}
{\True $1\ m/s^2$}
{$1,5\ m/s^2$}
{$2\ m/s^2$}	
\haidong{14}
\loigiai{
\begin{itemize}
\item Đổi: $18\ km/h = 5\ m/s;\ 54\ km/h = 15\ m/s$.
\item Áp dụng công thức độc lập thời gian: $v_t^2-v_0^2=2as \Rightarrow a=\dfrac{v_t^2-v_0^2}{2s}=\dfrac{15^2-5^2}{2.100}=1\ m/s^2$. 
\end{itemize}	
}
\end{vidu}

\begin{vidu}%[1P1-5.3K][Loai1]
Khi ô tô đang chạy với tốc độ $10\ m/s$ trên đoạn đường thẳng thì người lái xe tăng ga và ô tô chuyển động nhanh dần đều. Sau $20\ s$, ô tô đạt tốc độ $14\ m/s$. Tốc độ trung bình trên đoạn đường đó là
\choice
{\True $12\ m/s$}
{$9\ m/s$}
{$4\ m/s$}
{$14\ m/s$}
\haidong{16}
\loigiai{
Ta có: $14 = 10 + 20a \Rightarrow a = 0,2\ s \Rightarrow s=10.20+\dfrac{1}{2}.0,2.20^2=240\ m \Rightarrow v_{tb}=\dfrac{s}{t}=\dfrac{240}{20}=12\ m/s$. 
}
\end{vidu}
\begin{vidu} %[1P1B5-3][Loai1] 
Một đoàn tàu rời ga chuyển động thẳng nhanh dần đều. Sau $1$ phút đạt tốc độ $40\ km/h$. Quãng đường mà tàu đi được trong $1$ phút đó là
\choice
{\True $333\ m$}
{$666\ m$}
{$20\ m$}
{$240\ m$}
\haidong{15}
\loigiai{		
\begin{itemize}
\item Đổi $40\  km/h = 11,1\  m/s$.
\item Gia tốc của tàu $a = \dfrac{v_t}{t} = \dfrac{11,1}{60} = 0,185\  m/s^2$.
\item Quãng đường mà tàu đi được trong 1 phút: $s = \dfrac{at^2}{2} = \dfrac{1}{2}0,185.60^2 = 333$ m.
\end{itemize}}
\end{vidu}
\loai{Chuyển động chậm dần đều}
\begin{vidu}%[1P1-5.2K][Loai1]
Một chất điểm đang chuyển động với tốc độ $12\ km/h$ thì chuyển động thẳng chậm dần đều, sau $1$ phút thì dừng lại. Gia tốc của chất điểm có độ lớn là
\choice
{\True $0,056\ m/s^2$}
{$200\ m/s^2$}
{$0,56\ m/s^2$}
{$2\ m/s^2$}
\haidong{14}
\loigiai{
\begin{itemize}
\item $12\ km/h = \dfrac{10}{3}$ m/s; 1 phút $=$ 60 s.
\item Ta có: $v =v_0+at \Rightarrow a=\dfrac{v-v_0}{t}=\dfrac{0-\dfrac{10}{3}}{60}=-\dfrac{1}{18}=-0,056\ m/s^2 \Rightarrow \left|a\right| =0,056\ m/s^2$. 
\end{itemize}	
}
\end{vidu}

\begin{vidu}%[1P1-5.2K][Loai1]
Một xe đang chuyển động với tốc độ $36\ km/h$ thì hãm phanh, chuyển động chậm dần đều, sau $20$ giây tốc độ là $18\ km/h$. Tốc độ sau khi hãm được $30\ s$ là 
\choice
{$4\ m/s$}
{$3\ m/s$}
{\True $2,5\ m/s$}
{$1\ m/s$}
\haidong{15}
\loigiai{
\begin{itemize}
\item Đổi $36\ km/h = 10\ m/s;\ 18\ km/h = 5\ m/s$.
\item Ta có: $v =v_0+at \Rightarrow 5=10+20a \Rightarrow a=-0,25\ m/s^2$.  
\item Vận tốc của vật sau 30 s kể từ lúc hãm phanh là: $v=10-0,25.30=2,5\ m/s$. 
\end{itemize}	
}
%<MyLT>
\end{vidu}

\loai{Chuyển động trên nhiều giai đoạn}
\begin{vidu} %[1P1-5.3K][Loai1] 
Một đoàn tàu rời ga chuyển động với gia tốc $0,1\ m/s^2$ trên một đoạn đường dài $500\ m$, sau đó thì chuyển động đều. Sau $1$ giờ tàu đi được quãng đường là
\choice
{$35\ km$}
{\True $35,5\ km$}
{$40\ km$}
{$40,5\ km$}
\haidong{16}
\loigiai{
\begin{itemize}
\item Gọi $t_1$ là thời gian tàu rời ga với gia tốc 0,1 $ m/s^2$.
\item Ta có: $t_1 = \sqrt {\dfrac{2s}{a}} = \sqrt{\dfrac{2.500}{0,1}} = 100\ s.$
\item Vận tốc v của tàu sau thời gian $t_1$ tính bởi $v = at_1 = 0,1.100 = 10$  m/s.
\item Thời gian $t_2$ tàu chạy đều với vận tốc v là $t_2 = 1h - t_1 = 3600 - 100 = 3500$ s.
\item Đoạn đường tàu chạy đều $s_2 = v.t_2 = 10.3500 = 35 000$ m.
\item Đoạn đường tàu chạy được sau 1h là: $s = s_1 + s_2 = 500 + 35000 = 35500\ m = 35,5\ km$.
\end{itemize}	}
\end{vidu}

\begin{vidu}%[1P1T5-3][Loai1] 
Một xe đang chuyển động với tốc độ $7,2\ km/h$ thì tăng tốc. Sau $4\ s$, xe đi thêm được $40\ m$. Cuối giây thứ $6$, xe tắt máy, sau $13\ s$ thì ngừng hẳn lại. Quãng đường xe đi thêm được kể từ khi tắt máy là
\choice
{\True $169\ m$}
{$100\ m$}
{$141\ m$}
{$123\ m$}
\haidong{17}
\loigiai{
\begin{itemize}
\item Ta có: $s=v_0t+\dfrac{a}{2}t^2\Leftrightarrow 40=2.4+\dfrac{a}{2}.4^2\Rightarrow a=4\ m/s^2$.
\item Lại có:  $v_t=v_0+at=2+4.6=26\ m/s$.
\item Vận tốc của vật lúc $t = 6$  s: $v_t=v_0+at\Rightarrow v_t=2+4.6=26\ m/s$.
\item Chọn gốc thời gian lúc xe lúc xe bắt đầu tắt máy. 
\item Sau khi xe tắt máy thì nó chuyển động chậm dần đều với vận tốc đầu $v_0'=26\ m/s$ cho đến khi vận tốc bằng 0 thì hết thời gian là $t = 13$  s. 
\item Do đó:
\begin{align*}
v_t'=v_0'+a't\Leftrightarrow 0=26+13a'\Rightarrow a'=-2\ m/s^2.
\end{align*}
\item Quãng đường đi trong 13 giây cuối là: $s_{(t=13s)}=v_0't+\dfrac{1}{2}a't^2=169\ m$.
\end{itemize}
\begin{note}
Không đươc áp dụng cách tính lấy quãng đường trong 19 s trừ đi quãng đường trong 6 s để được quãng đường trong 13 s cuối. Vì có hai quá trình chuyển động (nhanh dần đều và chậm dần đều).
\end{note}	
}
\end{vidu}
\loai{Khai thác số liệu từ phương trình}
\begin{vidu}%[1P1K5-2][Loai1]
Phương trình độ dịch chuyển của một chất điểm có dạng: $d = 10t + 4t^2$ $(x:\ m;\ t:\ s)$. Lúc $t = 2\ s$ vận tốc của vật là
\choice
{$28\ m/s$}
{$18\ m/s$}
{$16\ m/s$}
{\True $26\ m/s$}	
\haidong{14}
\loigiai{
\begin{itemize}
\item Ta có: $d=v_0t+\dfrac{at^2}{2} \Rightarrow d=10t+4t^2=10t+\dfrac{8t^2}{2}$. 
\item Vận tốc ban đầu của vật là: $v_0=10\ m/s$. 
\item Gia tốc của vật là $a = 8\ m/s^2$.
\item Phương trình vận tốc của vật là: $v_t = 10 + 8t$ m/s.
\item Vận tốc của vật lúc $t = 2$ s là $v_t = 10 + 8.2 = 26$ m/s.
\end{itemize}	
}
\end{vidu}
%\newpage

\loai{Vận dụng công thức độc lập thời gian}
\begin{vidu}%[1P1K5-2][Loai1]
Một đoàn tàu rời ga, chuyển động nhanh dần đều từ trạng thái nghỉ, sau khi đi được $1000\ m$ thì đạt đến tốc độ $10\ m/s$. Tốc độ của tàu sau khi đi được $2000\ m$ là
\choice
{\True $10\sqrt{2}\ m/s$}
{$20\sqrt{2}\ m/s$}
{$10\sqrt{3}\ m/s$}
{$10\sqrt{5}\ m/s$}
\loigiai{
Ta có: $\begin{cases}
v_1^2-v_0^2=2as_1\Rightarrow a=\dfrac{v_1^2-v_0^2}{2s_1}=\dfrac{10^2-0}{2.1000}=0,05\ m/s^2\\
v_2^2-v_0^2=2as_2\Rightarrow v=\sqrt{2as_2+v_0^2}=10\sqrt{2}\ m/s.
\end{cases}$
}
\haidong{10}
\end{vidu}

\loai{Quãng đường vật đi trong giây thứ n}

\begin{vidu}%[1P1-5.3K][Loai1] 
Một ô tô đang chuyển động với tốc độ $12\ m/s$ thì hãm phanh, chuyển động chậm dần đều và đi thêm $36\ m$ thì dừng lại. Quãng đường ô tô đi được trong $2$ giây cuối cùng trước khi dừng hẳn là
\choice
{\True $4\ m$}
{$5\ m$}
{$6\ m$}
{$7\ m$}
\haidong{10}
\loigiai{
\begin{itemize}
\item $v^2-v_0^2=2as\Rightarrow a=\dfrac{-v_0^2}{2s}=-2\ m/s^2\Rightarrow s=v_0t+\dfrac{a}{2}t^2\Leftrightarrow 36=12t-t^2\Rightarrow t=6\ s$.
\item $s=v_0t+\dfrac{a}{2}t^2\Rightarrow \Delta s=s_{(t=6)}-s_{(t=4)}=12\left(6-4\right)-\dfrac{2}{2}\left(6^2-4^2\right)=4\ m$.
\end{itemize}	
}
\end{vidu}
\begin{vidu}%[1P1G5-3][Loai1] 
Một vật chuyển động nhanh dần đều từ trạng thái đứng yên. Trong giây thứ $4$ đi được $7\ m$. Quãng đường đi trong giây thứ $5$ là
\choice
{\True $9\ m$}
{$6\ m$}
{$7\ m$}
{$8\ m$}	
\haidong{12}
\loigiai{
\begin{itemize}
\item Vì vật bắt đầu chuyển động từ trạng thái đứng yên nên $v_0 = 0$. 
\item Quãng đường vật đi được trong 4 s đầu: $s_4=\dfrac{1}{2}at_4^2=8a$.
\item Quãng đường vật đi được trong 3 s đầu: $s_3=\dfrac{1}{2}at_3^2=4,5a$.
\item Theo đề ra: $s_4-s_3=7\Leftrightarrow 7=8a-4,5a\Rightarrow a=2\ m/s^2$.
\item Quãng đường vật đi được trong 5 s đầu: $s_5=\dfrac{1}{2}at_5^2=25\ m$.
\item Quãng đường vật đi được trong 4 s đầu: $s_4=\dfrac{1}{2}at_4^2=16\ m$.
\item Quãng đường đi được trong giây thứ 5: $\Delta s=s_5-s_4=9\ m$.
\end{itemize}
}
\end{vidu}
\loai{Thời gian vật đi mét thứ n}
\begin{vidu}%[1P1-5.3K][Loai1] 
Một vật bắt đầu chuyển động nhanh dần đều từ trạng thái đứng yên và đi hết quãng đường trong thời gian $t = 2\ s$. Thời gian để vật đi được nửa cuối của đường đi là
\choice
{$2+\sqrt{2}\ s$}
{$1+\sqrt{2}\ s$}
{\True $2-\sqrt{2}\ s$}
{$3-\sqrt{2}\ s$}	
\haidong{14}
\loigiai{
\begin{itemize}
\item Khi đi cả đoạn đường s, mất thời gian t nên ta có: $s=\dfrac{1}{2}at^2\quad (1)$
\item Khi đi nửa đoạn đường đầu mất thời gian $t_1$ nên: $\dfrac{s}{2}=\dfrac{1}{2}at_1^2\quad (2)$
\item Từ $(1)$ và $(2)$ ta có: $\dfrac{t^2}{t_1^2}=2\Rightarrow t_1=\dfrac{t\sqrt{2}}{2}$.
\item Vậy thời gian trong nửa cuối của quãng đường là: $\Delta t=t-t_1=2-\dfrac{t\sqrt{2}}{2}=2-\sqrt{2}$.
\end{itemize}
}
\end{vidu}
\begin{vidu}%[1P1T5-3][Loai1] 
Một viên bi chuyển động thẳng nhanh dần đều không vận tốc đầu, xuất phát trên đỉnh một máng nghiêng dài $10\ m$ và trong giây thứ $5$ nó đi được quãng đường bằng $36\ cm$. Thời gian để vật đi hết $1$ mét cuối cùng trên máng nghiêng là
\choice
{$0,51\ s$}
{\True $0,81\ s$}
{$0,33\ s$}
{$1,56\ s$}
\haidong{14}
\loigiai{
\begin{itemize}
\item Quãng đường vật đi được sau 4s và sau 5 s đầu tiên là: $\begin{cases}
s_4=\dfrac{1}{2}a.4^2=8a.\\
s_5=\dfrac{1}{2}a.5^2=12,5a.
\end{cases}$
\item Quãng đường bi đi được trong giây thứ 5 là: 	$\Delta s=s_5-s_4\Leftrightarrow 36=4,5a\Rightarrow a=8\ cm/s^2$.
\item Gọi thời gian để vật đi hết 9 m đầu và 10 m đầu là $t_9, t_{10}$ ta có: 	$\begin{cases}
9=\dfrac{1}{2}at_9^2.\\
10=\dfrac{1}{2}at_10^2.\\
\end{cases}\Rightarrow \begin{cases}
t_9=\sqrt{\dfrac{18}{a}}.\\
t_{10}=\sqrt{\dfrac{20}{a}}.\\
\end{cases}$
\item Thời gian để vật đi hết 1 m cuối là: $\Delta t=t_{10}-t_9=\sqrt{\dfrac{20}{0,08}}-\sqrt{\dfrac{18}{0,08}}=0,81\ s$.
\end{itemize}
}
\end{vidu}











%mac
\loai{Nhận biết đồ thị}
\begin{vidu}%SBT9.3
Đồ thị nào sau đây là của chuyển động thẳng chậm dần đều?
\begin{center}
\begin{tikzpicture}[scale=1,thick,>=stealth',x=1.5cm,y=0.2cm,draw=Mapcolor]
\tkzDefPoints{0/0/o, 0/12/x, 2/0/t, 0/8/m, 1.5/8/n, 2/0/n1, 0/12/n2}
\tkzDrawLines[very thick,color=Mapcolor,add = 0 and 0](m,n)
\tkzDrawSegments[->](o,x o,t)
\tkzLabelPoint[below](o){$O$}
\tkzLabelPoint[right](x){$v$}
\tkzLabelPoint[above](t){$t$}
\node at (1,-1.5){(a)};
\begin{scope}[shift={(3,0)}]
\tkzDefPoints{0/0/o, 0/12/x, 2/0/t, 0/8/m, 1.5/11/n, 2/0/n1, 0/12/n2}
\tkzDrawLines[very thick,color=Mapcolor,add = 0 and 0](m,n)
\tkzDrawSegments[->](o,x o,t)
\tkzLabelPoint[below](o){$O$}
\tkzLabelPoint[right](x){$v$}
\tkzLabelPoint[above](t){$t$}
\node at (1,-1.5){(b)};
\end{scope}
\begin{scope}[shift={(6,0)}]
\tkzDefPoints{0/0/o, 0/12/x, 2/0/t, 0/8/m, 1.5/0/n, 2/0/n1, 0/12/n2}
\tkzDrawLines[very thick,color=Mapcolor,add = 0 and 0](m,n)
\tkzDrawSegments[->](o,x o,t)
\tkzLabelPoint[below](o){$O$}
\tkzLabelPoint[right](x){$v$}
\tkzLabelPoint[above](t){$t$}
\node at (1,-1.5){(c)};
\end{scope}
\begin{scope}[shift={(9,0)}]
\tkzDefPoints{0/0/o, 0/12/x, 2/0/t, 0.2/8/m, 1.5/1/n, 2/0/n1, 0/12/n2}
\tkzDrawLines[very thick,color=Mapcolor,add = 0 and 0,bend right=30pt](m,n)
\tkzDrawSegments[->](o,x o,t)
\tkzLabelPoint[below](o){$O$}
\tkzLabelPoint[right](x){$v$}
\tkzLabelPoint[above](t){$t$}
\node at (1,-1.5){(d)};
\end{scope}
\end{tikzpicture}
\end{center}
\choice
{$a$}
{$b$}
{\True $c$}
{$d$}
\haidong{10} \end{vidu}
\loai{Từ đồ thị, nhận xét tính chất chuyển động}
\begin{vidu}
\immini[thm]{Dựa vào đồ thị $(v-t)$ của vật chuyển động trong hình bên. Mô tả chuyển động nào sau đây là đúng?
\choice
{Từ $0\ s$ - $40\ s$ vật chuyển động thẳng đều}
{\True Từ $40\ s - 80\ s$ vật chuyển động thẳng đều}
{Từ $0\ s-40\ s$ vật chuyển động thẳng chậm dần đều}
{Từ $40\ s-80\ s$ vật chuyển động thẳng biến đổi đều}}{
\begin{tikzpicture}[very thick,draw=Mapcolor]
\tkzDefPoints{0/0/o, 0/3.5/v, 4.5/0/t, 0/1/40a, 0/3/120, 1/0/40b, 2/0/80, 4/0/160, 1/3/x, 2/3/y, 4/0/z}
\tkzDrawSegments[->, thick](o,v o,t)
\tkzDrawSegments[Mapcolor,very thick](40a,x x,y y,z)
\tkzDrawSegments[dashed](120,x x,40b y,80)
\tkzLabelPoint[left](120){$120$}
\tkzLabelPoint[left](40a){$40$}
\tkzLabelPoint[below](40b){$40$}
\tkzLabelPoint[below](80){$80$}
\tkzLabelPoint[below](160){$160$}
\tkzLabelPoint[below](o){$O$}
\tkzLabelPoint[above](v){$v\ (m/s)$}
\tkzLabelPoint[above right](t){$t\ (s)$}
\end{tikzpicture}
}
\loigiai{
}
\haidong{10} \end{vidu}

\loai{Tính toán các đại lượng trên đồ thị 1 giai đoạn}
\begin{vidu}
\immini[thm]{Một chất điểm chuyển động thẳng biến đổi đều có đồ thị vận tốc theo thời gian như hình vẽ. Độ lớn gia tốc của chất điểm là
\choice
{$0,2\ m/s^2$}
{$0,4\ m/s^2$}
{\True $0,5\ m/s^2$}
{$0,6\ m/s^2$}}{\begin{tikzpicture}[very thick, scale=0.9,draw=Mapcolor]
\tkzDefPoints{0/0/x, 0/4.5/v, 5.5/0/t, 0/1.5/m, 5/4/n}
\tkzDrawSegments[->](x,t x,v)
\tkzDrawSegments[Mapcolor, very thick](m,n)
\draw  [help  lines, xstep=0.5cm,ystep=0.5cm]  (0,0)  grid  (5,4);

\tkzLabelPoint[left](m){$v_0$}
\tkzLabelPoint[above](v){$v ~(m/s)$}
\tkzLabelPoint[below right](t){$t~(s)$}
\tkzLabelPoint[below left](x){$O$}
\node[left] at (0,1) {2};
\node[left] at (0,2) {4};
\node[left] at (0,3) {6};
\node[left] at (0,4) {8};
\node[below] at (1,0) {2};
\node[below] at (2,0) {4};
\node[below] at (3,0) {6};
\node[below] at (4,0) {8};
\node[below] at (5,0) {10};
\end{tikzpicture}}
\haidong{2}

\loigiai{
}
\haidong{10} \end{vidu}
\begin{vidu}
\immini[thm]{Một chất điểm chuyển động thẳng có đồ thị như hình. Tính độ dịch chuyển của chất điểm trong thời gian từ lúc $t_1=4\ h$ đến lúc $t_2=7\ h$
\choice
{$90\ km$}
{\True $45\ km$}
{$165\ km$}
{$30\ km$}}{\begin{tikzpicture}[very thick, scale=1,draw=Mapcolor]
\tkzDefPoints{0/0/o, 0/3/v, 4/0/t}
\tkzDrawSegments[->, thick](o,v o,t)
\draw [very thick,Mapcolor](0,2.5) -- (2.5,0);
\draw [dashed, thin] (0,1) -- (1.5,1) -- (1.5,0);
\tkzLabelPoint[below](o){$O$}
\tkzLabelPoint[above](v){$v~(km/h)$}
\tkzLabelPoint[below](t){$t~(h)$}
\node[left] at (0,1) {30};
\node [below] at (1.5,0) {4};
\node[below] at (2.47,0) {7};
\end{tikzpicture}}
\haidong{4}
\loigiai{
}
\haidong{10} \end{vidu}
\loai{Tính toán các đại lượng trên đồ thị nhiều giai đoạn}
\begin{vidu}%[1P1-6.2K][Loai1]
\immini[thm]{Một chất điểm chuyển động thẳng có đồ thị vận tốc - thời gian như hình vẽ.
\begin{itemize}
\item[a)] Xác định loại chuyển động ứng với mỗi đoạn của đồ thị.
\item[b)] Xác định gia tốc tương ứng. 
\item[c)] Lập phương trình vận tốc ứng với từng đoạn trên đồ thị. 
\item[d)] Tính quãng đường vật đã đi.
\end{itemize}}
{\begin{tikzpicture}[scale=.9,thick,>=stealth',x=0.05cm,y=0.04cm,draw=Mapcolor]
\tkzDefPoints{0/0/o, 0/80/x, 100/0/t, 0/20/a, 20/60/b, 20/0/b1, 0/60/b2, 40/60/c, 40/0/c1, 0/60/c2, 80/0/d}
\tkzDrawSegments[dashed](b1,b b2,b c1,c)
\tkzDrawSegments[very thick, Mapcolor](a,b b,c c,d)
\tkzDrawSegments[->](o,x o,t)
\tkzLabelPoint[below left](o){$O$}
\tkzLabelPoint[left](b2){$60$}
\tkzLabelPoint[left](a){$20$}
\tkzLabelPoint[below](b1){$20$}
\tkzLabelPoint[below](c1){$40$}
\tkzLabelPoint[below](d){$80$}
\tkzLabelPoint[below right](a){$A$}
\tkzLabelPoint[above](b){$B$}
\tkzLabelPoint[above](c){$C$}
\tkzLabelPoint[above right](d){$D$}
\tkzLabelPoint[right](x){$v\ m/s$}
\tkzLabelPoint[below](t){$t\ (s)$}
\end{tikzpicture}}
\haidong{15}
\loigiai{
* Giai đoạn AB ($0 \le  t \le  20s$):
\begin{itemize}
\item Đồ thị dốc lên, do đó trong giai đoạn này chất điểm chuyển động nhanh dần đều với gia tốc được xác định theo công thức: $a_{AB}=\dfrac{v_B-v_A}{t_{AB}}=\dfrac{60-20}{20}=2\ m/s^2$
\item Phương trình vận tốc trên giai đoạn AB: $v_{AB}=v_A+a_{AB}t=20+2t\ \left(m/s;\ s\right)$
\item Quãng đường chất điểm đi được trong giai đoạn AB:
\begin{align*}
s_{AB}=\dfrac{v_B^2-v_A^2}{2a_{AB}}=\dfrac{60^2-20^2}{2.2}=800\ m.
\end{align*}
\end{itemize}
* Giai đoạn BC ($20\ s \le  t \le  40\ s$):
\begin{itemize}
\item Đồ thị nằm ngang, do đó trong giai đoạn này chất điểm chuyển động thẳng đều với gia tốc: $a_{BC}=0$.
\item Phương trình vận tốc trên giai đoạn BC: $v_{BC}=v_B=60\ m/s$.
\item Quãng đường chất điểm đi được trong giai đoạn BC:
$s_{BC}=v_{BC}.t_{BC}=60.20=1200\ m$.
\end{itemize}
* Giai đoạn CD ($40\ s \le  t \le  80\ s$):
\begin{itemize}
\item Đồ thị dốc xuống, do đó trong giai đoạn này chất điểm chuyển động chậm dần đều với gia tốc được xác định theo công thức: $a_{CD}=\dfrac{v_D-v_C}{t_{CD}}=\dfrac{0-60}{40}=-1,5\ m/s^2$.
\item Phương trình vận tốc trên CD: $v_{CD}=v_C+a_{CD}\left(t-t_0\right)=60-1,5\left(t-40\right)$
(v đo bằng m/s và t đo bằng s).
\item Quãng đường chất điểm đi được trong giai đoạn CD:
$s_{CD}=\dfrac{v_D^2-v_C^2}{2a_{CD}}=\dfrac{0^2-60^2}{2.\left(-1,5\right)}=1200\ m$.
\item Vậy tổng quãng đường khi đi từ A đến D là:
$s=s_{AB}+s_{BC}+s_{CD}=800+1200+1200=3200\ m$.
\end{itemize}}
\haidong{10} \end{vidu}

\begin{vidu}
\immini[thm]{Dựa vào mối quan hệ giữa gia tốc và độ dốc của đồ thị vận tốc - thời gian, hãy tính tỉ số về độ lớn gia tốc của vật trong khoảng thời gian $AB$ và $OA$ ở hình bên.
\choice
{$\dfrac{1}{3}$}
{$\dfrac{1}{2}$}
{$1$}
{\True $3$}}{\begin{tikzpicture}[very thick,scale=0.9,draw=Mapcolor]
\tkzDefPoints{0/0/o, 0/3/v, 5.5/0/t, 0/2/o', 3.5/2/a', 3.5/0/a, 4.5/0/b}
\tkzDrawSegments[->](o,t o,v)
\tkzDrawSegments[dashed](o',a' a',a)
\tkzDrawSegments[Mapcolor, very thick](o,a' a',b)
\tkzLabelPoint[above](v){$v ~(m/s)$}
\tkzLabelPoint[below right](t){$t~(s)$}
\tkzLabelPoint[below left](o){$O$}
\tkzLabelPoint[below](a){$A$}
\tkzLabelPoint[below](b){$B$}
\tkzLabelPoint[left](o'){$40$}
\tkzMarkAngles[size=0.6cm](a,o,a')
\tkzMarkAngles[size=0.4cm](a',b,a)
\tkzLabelAngles[pos=1](a,o,a'){$30^\circ$}
\tkzLabelAngles[pos=0.7](a',b,a){$60^\circ$}
\end{tikzpicture}}
\haidong{5}

\loigiai{
}
\haidong{10} \end{vidu}



